\documentclass[12pt]{article}
% єдина преамбула -- nmtbase.sty у корені проєкту (знаходиться і з кореня, і з теки теми, і з підтеки)
\IfFileExists{nmtbase.sty}{\usepackage{nmtbase}}{\IfFileExists{../nmtbase.sty}{\usepackage{../nmtbase}}{\usepackage{../../nmtbase}}}
\fancyhead[C]{\small\color{gray!80} База завдань НМТ 2022--2026 \textendash{} усі теми}
\fancyhead[R]{\small\color{mainGreen}\textbf{\href{https://t.me/pvtr2525}{@pvtr2525}}}
\fancyfoot[L]{\small\color{gray!80}\href{https://t.me/pvtr2525}{ТГ: @pvtr2525}}
\fancyfoot[C]{\small\color{gray!80}\thepage}
\fancyfoot[R]{\small\color{gray!80}НМТ 2022--2026}
\usepackage{docmute}
% у зведених документах: розділ = section, тема = subsection, рік = subsubsection
\setcounter{tocdepth}{3}
\renewcommand{\chapterTitle}[1]{\clearpage\phantomsection\addcontentsline{toc}{subsection}{#1}\sectionTitle{#1}}
\renewcommand{\typeTitle}[1]{\needspace{6\baselineskip}\phantomsection\addcontentsline{toc}{subsubsection}{\quad #1}%
\par\vspace{0.3cm}\noindent{\large\bfseries\color{yearOrange}#1}\par\vspace{0.08cm}%
\noindent{\color{yearOrange}\rule{\linewidth}{0.4pt}}\par\nopagebreak\vspace{0.2cm}}
\newcommand{\partTitle}[1]{\clearpage\phantomsection\addcontentsline{toc}{section}{#1}%
\vspace*{0.5cm}\begin{tcolorbox}[colback=titleBg, colframe=mainGreen!60, boxrule=0.8pt, arc=8pt, halign=center, fontupper=\LARGE\bfseries\color{mainGreen}]
#1
\end{tcolorbox}\par\vspace{0.4cm}}
\begin{document}
\thispagestyle{empty}
\vspace*{2cm}
\begin{center}
{\Huge\bfseries\color{mainGreen} БАЗА ЗАВДАНЬ НМТ}\\[0.4cm]
{\LARGE\bfseries\color{mainGreen} 2022--2026}\\[0.6cm]
{\Large\bfseries з математики}\\[0.8cm]
{\large Усі завдання основних сесій: 6 розділів, 43 теми}\\[1.2cm]
{\large\color{mainGreen}\textbf{\href{https://t.me/pvtr2525}{Телеграм: @pvtr2525}}}
\end{center}
\clearpage
{\hypersetup{linkcolor=black}\tableofcontents}
\partTitle{РОЗДІЛ 1. ЧИСЛА І ВИРАЗИ}
\input{"1.\space Числа.\space Звичайні\space та\space десяткові\space дроби.\space Модуль\space числа/завдання.tex"}
\input{"2.відсотки,\space арифметичні\space задачі,\space подільність/завдання.tex"}
\input{"3.\space Степені.\space Одночлени.\space стандартний\space вигляд\space числа/завдання.tex"}
\input{"4.Многочлени.\space ФСМ/завдання.tex"}
\input{"5.\space дробові\space раціональні\space вирази\space і\space перетворення/завдання.tex"}
\input{"7.\space Квадратні\space корені.\space Корені\space вищих\space степенів.\space Раціональний\space степінь/завдання.tex"}
\input{"26.\space Тригонометричні\space вирази/завдання.tex"}
\input{"30.\space Логарифм.\space Логарифмічна\space функція/завдання\space вирази.tex"}
\partTitle{РОЗДІЛ 2. РІВНЯННЯ І НЕРІВНОСТІ}
\input{"6.\space Лінійні\space рівняння/завдання.tex"}
\input{"8.\space квадратні\space рівняння,\space і\space рівняння,\space що\space до\space них\space зводяться/завдання.tex"}
\input{"9.\space Системи\space рівнянь/завдання.tex"}
\input{"19.\space Лінійні\space та\space квадратні\space нерівності.\space Системи\space лінійних\space нерівностей.\space Метод\space інтервалів/завдання.tex"}
\input{"22.\space Модульні\space \space рівняння\space та\space нерівності/завдання.tex"}
\input{"23.\space Ірраціональні\space рівняння\space та\space нерівності/завдання.tex"}
\input{"27.\space Тригонометричні\space рівняння/завдання.tex"}
\input{"28.\space Показникова\space функція.\space Показникові\space рівняння/Показникові\space рівняння/завдання.tex"}
\input{"29.\space Показникові\space нерівності/завдання.tex"}
\input{"31.\space Логарифмічні\space рівняння/завдання.tex"}
\input{"32.\space Логарифмічні\space нерівності/завдання.tex"}
\input{"43.\space Параметри/завдання.tex"}
\partTitle{РОЗДІЛ 3. ФУНКЦІЇ}
\input{"20.\space Функція\space та\space її\space властивості/завдання.tex"}
\input{"21.\space \space Графіки\space елементарних\space функцій,\space їх\space побудова\space графіків\space функцій\space та\space їх\space перетворення/завдання.tex"}
\input{"24.\space Арифметична\space прогресія/завдання.tex"}
\input{"25.\space Геометрична\space прогресія/завдання.tex"}
\input{"28.\space Показникова\space функція.\space Показникові\space рівняння/Показникові\space функція\space і\space вирази/завдання.tex"}
\input{"30.\space Логарифм.\space Логарифмічна\space функція/завдання\space функція.tex"}
\input{"33.\space Похідна/завдання.tex"}
\input{"34.\space Первісна/завдання.tex"}
\partTitle{РОЗДІЛ 4. ПЛАНІМЕТРІЯ}
\input{"10.\space Основи\space планіметрії/завдання.tex"}
\input{"11.\space Трикутники/завдання.tex"}
\input{"12.\space Паралелограм/завдання.tex"}
\input{"13.\space Прямокутник/завдання.tex"}
\input{"14.\space Квадрат/завдання.tex"}
\input{"15.\space Ромб/завдання.tex"}
\input{"16.\space Трапеція/завдання.tex"}
\input{"17.\space Коло,\space круг\space і\space їх\space елементи/завдання.tex"}
\partTitle{РОЗДІЛ 5. СТЕРЕОМЕТРІЯ}
\input{"18.\space Координати\space і\space вектори\space на\space площині\space і\space в\space просторі.\space Рівняння\space кола/завдання.tex"}
\input{"38.\space Призма.\space Паралелепіпед.\space Куб/завдання.tex"}
\input{"39.\space Піраміда/завдання.tex"}
\input{"40.\space Циліндр/завдання.tex"}
\input{"41.\space Конус/завдання.tex"}
\input{"42.\space Куля.\space Сфера/завдання.tex"}
\partTitle{РОЗДІЛ 6. КОМБІНАТОРИКА, ТЕОРІЯ ЙМОВІРНОСТІ, СТАТИСТИКА}
\input{"35.\space Комбінаторика/завдання.tex"}
\input{"36.\space Теорія\space ймовірності/завдання.tex"}
\input{"37.\space Статистика/завдання.tex"}
\end{document}
